\section{Total Cost of Ownership and Return on Investment}
\label{sec:tco}
% \Vzero/\Vone are defined in the main preamble; guard \Vtwo so this fold
% is self-contained. \providecommand does not clobber an existing def.
\providecommand{\Vtwo}{\textcolor{skyblue}{V2}}

A co-designed accelerator only pays for itself if the recurring efficiency it
buys eventually amortizes its one-time silicon cost. To evaluate this we adopt
FAST's total-cost-of-ownership (TCO) methodology~\citep{zhang_fast}
\emph{verbatim} (the same cost model, the same non-recurring-engineering (NRE)
decomposition, and the same return-on-investment (ROI) break-even analysis;
their Eq.~1--2 and Sec.~5.1), so that our deployment economics are read on
FAST's own terms and are directly comparable to their Table~4. We stress at the
outset that this is the weakest-grounded part of our evaluation: absolute TCO
depends on proprietary fleet data neither we nor FAST possess, several cost
constants are assumption-flagged, and, most importantly, the speedup that
drives the ROI is measured against a different baseline than FAST's
(\cref{par:tco_nonparity}). We therefore report ROI as an order-of-magnitude
plausibility check, not a precise claim, and we state every caveat inline.

\paragraph{FAST's cost model (what we mirror).}
FAST models the three-year cost of a fleet that offloads $n$ baseline
accelerator-units of traffic onto a co-designed chip as
%
\begin{equation}
\label{eq:tco_old}
\mathrm{TCO}_{\mathrm{old}}(n) \;=\; C_{\mathrm{cap}}(n) \;+\; t_D\, C_{\mathrm{op}}(n),
\end{equation}
%
capital cost plus lifetime ($t_D=3$~yr) operating cost, and defines the ROI of
committing the co-design NRE as
%
\begin{equation}
\label{eq:tco_roi}
\mathrm{ROI} \;=\; \frac{\mathrm{TCO}_{\mathrm{old}}\,(S-1)}
{\bigl(t_{\mathrm{design}}\,C_{\mathrm{eng}} + C_{\mathrm{mask}} + C_{\mathrm{IP}}\bigr)\,S},
\end{equation}
%
where $S$ is the per-unit efficiency gain of the co-designed chip over the
baseline; $\mathrm{ROI}>1$ means the deployment is profitable. Because per-chip
TCO is confidential, FAST uses $S=\text{Perf/TDP}$ as a proxy for Perf/TCO (the
two are reported to be highly correlated for datacenter
accelerators~\citep{jouppi2017datacenter,jouppi2021ten}) and anchors the
baseline unit to a commodity DGX-A100 accelerator. We reuse this proxy, this
DGX-anchored baseline unit, and FAST's engineering-NRE constant
($t_{\mathrm{design}}\,C_{\mathrm{eng}}=\text{65 engineer-years}\times\$240\mathrm{K}\times1.65$
overhead $=\$25.74\mathrm{M}$) unchanged. Inverting \cref{eq:tco_roi} at a
target ROI multiple $R$ gives the break-even deployment volume
%
\begin{equation}
\label{eq:tco_breakeven}
n^\star(R) \;=\; \frac{R\,\cdot\,\mathrm{NRE}\,\cdot\,S}{(S-1)\,\cdot\,\mathrm{TCO}_{\mathrm{unit}}},
\qquad
\mathrm{NRE} = t_{\mathrm{design}}\,C_{\mathrm{eng}} + C_{\mathrm{mask}} + C_{\mathrm{IP}},
\end{equation}
%
the number of baseline units whose traffic must be offloaded before the NRE is
repaid. All arithmetic below is a pure evaluation of
\cref{eq:tco_old,eq:tco_roi,eq:tco_breakeven} on the eight co-designed
blueprints of \cref{app:blueprints}; no solver is involved.

\paragraph{Break-even under FAST parity.}
Under FAST's Perf/TDP proxy, our co-designed chips retain
\todo{RECONCILE with \cref{sec:perf_tdp}, which reports $2.73$--$4.83\times$.
Those are SERVING designs under the FAST TDP-ratio convention; these are PREFILL
blueprints after normalizing away the MAC count. Different baselines and
different conventions. State this in both places or the paper reads as
self-contradictory six pages apart.}$S=1.15\text{--}1.79\times$ Perf/TDP over a modeled TPU-v3 after normalizing away
raw MAC scale (measured against a naively scheduled baseline; see
\cref{par:tco_nonparity}). For the best chip (the DeepSeek-V3~\citep{deepseekv3report}
$\Vone$ design on TSMC~N7 with LPDDR5X, $S=1.79$), \cref{eq:tco_breakeven} gives a
$1\times$-ROI break-even of $\mathbf{4{,}410}$ units (\cref{tab:tco_breakeven},
\cref{fig:tco_breakeven}). This sits at the \emph{same order of magnitude} as
FAST's own reported $2{,}164\text{--}3{,}534$ units (their Table~4, at
$S=1.84\text{--}3.91$), and well within the thousands-to-tens-of-thousands scale
of modern accelerator fleets. As \cref{eq:tco_breakeven} makes explicit and
\cref{fig:tco_breakeven} shows, deployment volume dominates ROI: the ranking of
chips by break-even volume is driven jointly by $S$ and by NRE, and the
N7-plus-commodity-memory family ($\Vone$) is uniformly ROI-preferred over the
N5-plus-HBM3 family ($\Vtwo$), which carries ${\sim}\$15\mathrm{M}$ more mask NRE,
$1.7\text{--}2.5\times$ the unit TCO, and (except on the 13B) \emph{lower}
Perf/TDP, pushing its break-even to $13{,}000\text{--}19{,}000$ units.

\begin{table}[t]
\centering
\small
\todo{FIGURE NOT RENDERED: the graphics include is commented out, so this prints as a caption with no plot, yet it is cited in the running text.}\caption{FAST-parity break-even volumes (convention F: $S=\text{Perf/TDP}$
proxy, DGX-A100-anchored baseline unit), computed from
\cref{eq:tco_breakeven}. NRE is $\$50.7\mathrm{M}$ for the N7 ($\Vone$) family
and $\$65.7\mathrm{M}$ for the N5+HBM3 ($\Vtwo$) family. For reference, FAST's
Table~4 reports $1\times$-ROI volumes of $2{,}164\text{--}3{,}534$ units at
$S=1.84\text{--}3.91$. Cost constants are assumption-flagged; see
\cref{par:tco_nre}.}
\label{tab:tco_breakeven}
\begin{tabular}{@{}llrrrr@{}}
\toprule
\textbf{Chip (blueprint)} & \textbf{Process / memory} & $\mathbf{S}$ &
\textbf{1$\times$} & \textbf{2$\times$} & \textbf{4$\times$} \\
\midrule
DeepSeek-V3 $\Vone$ & N7 / LPDDR5X & $1.79$ & $\mathbf{4{,}410}$ & $8{,}821$ & $17{,}641$ \\
Llama1-65B $\Vone$  & N7 / LPDDR5X & $1.74$ & $4{,}591$ & $9{,}182$ & $18{,}363$ \\
Llama3-70B $\Vone$  & N7 / LPDDR5X & $1.54$ & $5{,}528$ & $11{,}057$ & $22{,}114$ \\
Llama1-13B $\Vtwo$  & N5 / HBM3    & $1.57$ & $6{,}920$ & $13{,}841$ & $27{,}681$ \\
Llama1-13B $\Vone$  & N7 / GDDR6X  & $1.42$ & $6{,}641$ & $13{,}283$ & $26{,}566$ \\
DeepSeek-V3 $\Vtwo$ & N5 / HBM3    & $1.24$ & $13{,}083$ & $26{,}166$ & $52{,}332$ \\
Llama1-65B $\Vtwo$  & N5 / HBM3    & $1.20$ & $15{,}079$ & $30{,}158$ & $60{,}317$ \\
Llama3-70B $\Vtwo$  & N5 / HBM3    & $1.15$ & $19{,}237$ & $38{,}473$ & $76{,}947$ \\
\bottomrule
\end{tabular}
\end{table}

\begin{figure}[t]
\centering
% TODO: render from tco_data.json (fast_convention/tco_data.json,
%       breakeven_volume_conventionF, 1x_ROI per representative_chips).
%       Grouped bar chart: one bar per blueprint (8), N7/V1 vs N5/V2
%       families in two colors; y-axis = 1x-ROI break-even volume (units,
%       log or linear). Overlay a horizontal shaded band at 2,164-3,534 =
%       FAST's Table-4 range, and a dashed marker at 3,650 = our best chip
%       under FAST's implied ~$42M NRE (\cref{par:tco_nre}).
% \includegraphics[width=0.78\linewidth]{figs/tco_breakeven.pdf}
\todo{FIGURE NOT RENDERED: the graphics include is commented out, so this prints as a caption with no plot, yet it is cited in the running text.}\caption{FAST-parity $1\times$-ROI break-even deployment volume per co-designed
blueprint (convention F), from \cref{eq:tco_breakeven}. The shaded band marks
FAST's own reported range ($2{,}164\text{--}3{,}534$ units); our best chip
(DeepSeek-V3 $\Vone$, $4{,}410$ units) lands just above it, and drops to
${\sim}3{,}650$ units (\emph{inside} FAST's range) when evaluated at FAST's
implied ${\sim}\$42\mathrm{M}$ total NRE. The N7$+$commodity-memory ($\Vone$)
family is uniformly ROI-preferred over N5$+$HBM3 ($\Vtwo$).}
\label{fig:tco_breakeven}
\end{figure}

\paragraph{Reconciling the NRE.}
\label{par:tco_nre}
The NRE terms in \cref{eq:tco_breakeven} carry the largest uncertainty, and we
flag them as assumptions~[A] rather than measured quantities: the mask-set cost
($\$15\mathrm{M}$ at N7, $\$30\mathrm{M}$ at N5), the IP/PHY-licensing cost
($\$10\mathrm{M}$, ASIC-Clouds-class~\citep{magaki2016asic}), the defect density
$D_0$ used in the die-yield model, and the PUE all follow public industry-class
estimates, not vendor data. FAST published only its engineering-NRE constant, so
rather than assert NRE parity we \emph{reconcile}: back-computing from FAST's own
Table~4 break-even volumes and reported $S$ implies a FAST total NRE of
${\sim}\$42.0\mathrm{M}$. Our best chip carries a larger, deliberately
conservative $\$50.7\mathrm{M}$; re-evaluating \cref{eq:tco_breakeven} at FAST's
implied ${\sim}\$42\mathrm{M}$ moves our best-chip break-even from $4{,}410$ to
$\mathbf{{\sim}3{,}650}$ units, now \emph{inside} FAST's own
$2{,}164\text{--}3{,}534$ range. The Perf/TDP proxy $S$ is independent of these
[A] constants except through NRE, so the headline break-even row is robust to
them: doubling the mask cost alone moves the $1\times$-ROI volume only to
${\sim}5{,}700$. We also disclose the honest error bar from the proxy itself: a
fully modeled TCO on both sides (convention M) inflates break-even to
${\sim}79{,}000\text{--}108{,}000$ units (an ${\sim}18\times$ spread), which is
precisely why FAST, and we, headline Perf/TDP and cite the Perf/TCO correlation
rather than claim TCO precision.

\paragraph{The zero-NRE deployment option.}
A distinct and strictly cleaner deployment story requires \emph{no silicon at
all}. Because CHEETAH's scheduler is a drop-in software layer, applying the joint
LP schedule to the \emph{existing} TPU-v3 fleet (same chip, same TDP) delivers
$S=2.14\text{--}2.51\times$ Perf/TDP on the LLM workloads ($2.21\times$ on
EfficientNet-B7~\citep{tan2019efficientnet}, $1.50\times$ on the Stable-Diffusion
U-Net) at software-only engineering cost, with $C_{\mathrm{mask}}=C_{\mathrm{IP}}=0$
(\cref{fig:tco_perftdp}). Under \cref{eq:tco_roi} the NRE collapses to a software
term and the ROI exceeds any silicon option at every volume. This is a genuine,
low-risk path to fleet-wide efficiency that does not depend on any of the
[A]-flagged silicon constants above.

\begin{figure}[t]
\centering
% TODO: render from tco_data.json (fast_convention/tco_data.json,
%       S_conventionF_perf_per_tdp_proxy for silicon chips;
%       zero_nre_option.finding for the TPU-v3 scheduling bars) and
%       fast_convention/table5_cheetah.md. Grouped bar chart:
%       (a) co-designed silicon 1.15-1.79x Perf/TDP, (b) zero-NRE TPU-v3
%       scheduling 2.14-2.51x (LLMs) / 2.21x (B7) / 1.50x (SD), against a
%       baseline of 1.0. Annotate the ~0.8x well-scheduled-TPU-v3 line as a
%       dashed threshold (\cref{par:tco_nonparity}).
% \includegraphics[width=0.78\linewidth]{figs/tco_perf_tdp.pdf}
\caption{Perf/TDP (FAST's TCO proxy) of the two deployment options against a
TPU-v3 baseline of $1.0$. Co-designed silicon retains $1.15\text{--}1.79\times$
after normalizing away MAC scale; deploying CHEETAH \emph{scheduling} on the
existing TPU-v3 fleet delivers $2.14\text{--}2.51\times$ at zero silicon NRE. The
dashed line marks the ${\sim}0.8\times$ figure obtained when the baseline is a
\emph{well-scheduled} TPU-v3 rather than a naive one: the baseline-non-parity
caveat of \cref{par:tco_nonparity}.}
\label{fig:tco_perftdp}
\end{figure}

\paragraph{Baseline non-parity: the load-bearing caveat.}
\label{par:tco_nonparity}
The single largest caveat on the silicon TCO case is that our $S$ and FAST's $S$
are \emph{not measured against the same baseline}. CHEETAH's speedup is measured
against a naive streaming schedule on the reference hardware, whereas FAST's is
measured against a production-calibrated simulator. This matters directly for
\cref{eq:tco_roi}: when the baseline is instead a \emph{well-scheduled} TPU-v3,
the co-designed chip's Perf/TDP is only ${\sim}0.8\times$ ($<1$), so
$S-1<0$ and \emph{no FAST-parity break-even exists at any deployment volume}. We
do not hide this. In that stricter accounting the silicon TCO case does not rest
on Perf/TDP parity at all; it rests on two things the scheduler alone cannot buy:
(i) the $34\text{--}40\times$ raw per-chip throughput and latency of the
co-designed datapath (and the corresponding reduction in fleet size and physical
footprint), mirroring FAST's own observation that latency and model-enablement
can justify even break-even ROI; and (ii) the zero-NRE scheduling option above,
whose ROI is unbounded because it commits no silicon NRE. Read this section
accordingly: the FAST-parity break-even numbers are a like-for-like comparison
under FAST's own convention, and the well-scheduled-baseline figure is the honest
floor; the two together, not either alone, are the deployment argument.
